% claim-ref: mosquitofish#dimensions
% claim-ref: mosquitofish#water-preparation
% claim-ref: mosquitofish#readiness
% claim-ref: mosquitofish#containment
% claim-ref: mosquitofish#equipment
% claim-ref: mosquitofish#stocking-gate
\CompanionHeader{TANK SETUP}{Outdoor habitat}{Western mosquitofish / unmeasured container}{Readiness must be established before stocking; no population count is prescribed.}
\Context{District depth recommendation: at least 24 inches. About 50 US gallons of actual water is a conservative planning proposal, NOT an official minimum or capacity guarantee. Ask the district to assess the actual dimensions. [MF-SM; proposed]}
\Panel{Measure the real system}{Record usable depth, footprint, actual volume after displacement, fish-safe material, stable base and sun/wind exposure. Provide steep sides, refuges and open water; prevent predator access and protect children/pets around deep water. [MF-SM; proposed]}
\Panel{Treat every refill}{Pre-treat every tap-water addition, including top-ups, for chloramine with a fish-appropriate product used as directed. Standing water is insufficient. Exclude runoff, pesticides, fertilizers, soap and spray drift. [MF-SM; MF-RETURN]}
\Panel{Prove the cycle without animals}{Establish biofiltration and document processing of a controlled waste load. Zero tests in unused clean water and elapsed weeks do not prove readiness. Require undetectable ammonia/nitrite on suitable tests; record units, pH, alkalinity and temperature. Do not chase chemistry with abrupt corrections. [MF-UF; MF-VET; proposed]}
\Panel{Contain adults and fry}{Protect intakes; screen overflow for fry with adequate flow and freeboard. Clean screens before storms: blockage can cause overtopping. No route to drains or natural waters. Keep gentle circulation/aeration, calm feeding areas and an open surface. [MF-FAO; MF-VET; proposed]}
\Panel{Season and equipment gate}{Log min/max water at fish depth, especially dawn, hot afternoons and cold spells. Prepare a protected contingency before freezing; avoid abrupt transfers. If heat is needed, use only manufacturer-approved outdoor equipment. Verify GFCI, weatherproof connections and all instructions before operation. [MF-SM; MF-CPSC; MF-VET; proposed]}
\Panel{Readiness record / all fields unmeasured}{Volume / depth / footprint / material: \hrulefill\par
Water min/max / pH / alkalinity: \hrulefill\par
Ammonia / nitrite / nitrate, units and cycle evidence: \hrulefill\par
Treatment / overflow / power / contingency: \hrulefill\par
Supplier identity / district review / capacity / permission: \hrulefill}
\Note{Before fish arrive}{Confirm every readiness field and district advice on initial count and acclimation. This page supplies neither a safe maximum population nor a universal optimum water range.}
\Sources{Clickable evidence: \href{https://www.smcmvcd.org/mosquitofish}{MF-SM}; \href{https://www.smcmvcd.org/request-mosquitofish-for-your-backyard-pond-or-water-garden}{MF-RETURN}; \href{https://ask.ifas.ufl.edu/publication/FA031}{MF-UF}; \href{https://www.msdvetmanual.com/exotic-and-laboratory-animals/aquarium-fish/management-of-aquarium-fish}{MF-VET}; \href{https://www.fao.org/fishery/docs/CDrom/FAO_Training/FAO_Training/General/t0581e/t0581e07.htm}{MF-FAO}; \href{https://www.cpsc.gov/s3fs-public/513.pdf}{MF-CPSC}. Source passages, dates and applicability: adjacent YAML. Planning only; actual habitat unverified.}{mosquitofish / tank-setup}
